\begin{conf-abstract}
{Investigating gasgeochemical production during the preparatory phase of dynamic brittle failure in Rotondo granite in the laboratory}
{Paul Selvadurai, Claudio Madonna, Clement Roques, Hanna Ventura, Matthias Brennwald, Rolf Kipfler, Benoit, Valley, Alexandra Lightfoot, Laurent Marguet}
{ETH Zurich}
{Laboratory geochemistry studies enhance our understanding of earthquake precursors. Controlled experimental conditions allow researchers to observe how microcracks, fluid flow, and mineral reactions influence geochemical signals, including gas emissions and isotope distributions. These experiments can identify specific geochemical changes, such as the release of radon or noble gases, associated with rock fracturing. Such insights are crucial for developing indicators of large events while calibrating geochemical signals, establishing links between these signals, stress states, and damage levels within rocks.

This study investigates the use of helium, argon, and CO2 as real-time geochemical tracers for brittle deformation in crystalline rocks, focusing on Rotondo granite from the Bedretto Laboratory. These gases, stored in mineral lattices or fluid inclusions, are released during microstructural damage; however, their behavior during progressive rock failure remains poorly understood. Uniaxial compression experiments monitored gas variability with a miniRuedi quadrupole mass spectrometer, complemented by spatial strain measurements using Distributed Fibre Optic Strain (DFOS) sensing, acoustic emission (AE) monitoring, and post-mortem X-ray computed tomography (XRCT) of fracture networks. The experimental design aims to identify distinct deformation stages, from crack initiation to macroscopic failure, and link these to transient noble gas signals. Preliminary results reveal stage-specific fluctuations in helium and CO2, supporting the hypothesis that gas emissions reflect structural damage evolution. This research advances the integration of mechanical and geochemical monitoring in rock mechanics, with important implications for underground stability assessment and early warning systems in geotechnical and seismic environments.}
\end{conf-abstract}
